% !PRINTED 223
\weberpage{223}

%<<CONT
also $\mathrm{sn}\,\Omega$ irgend eine Wurzel der eigentlichen Teilungsgleichung, so
hat die auf alle $\alpha$ auszudehnende Summe
\weq{13}{\Sigma(\pm 1)^\varepsilon(\pm 1)^{\varepsilon'}\ldots\, \mathrm{sn}\,\alpha\Omega = \psi(\Omega),}
worin die $k$ Vorzeichen von $\pm 1$ beliebig gewählt werden können, die Eigenschaft,
daß für jedes nach (12) bestimmte $\alpha$
\weq{14}{\psi(\alpha\Omega) = (\pm 1)^\varepsilon(\pm 1)^{\varepsilon'}\ldots \psi(\Omega)}
ist, und das Quadrat von $\psi(\Omega)$ bleibt durch die Substitutionen der Gruppe
$\mathfrak{B}_0$ ungeändert, ist also durch $n$te Einheitswurzeln und die Wurzeln
einer Transformationsgleichung rational ausdrückbar. $\psi(\Omega)$ ist also die
Quadratwurzel $\sqrt{A}$ aus einem solchen Ausdruck $A$. Die Anzahl der Ausdrücke (13)
beträgt aber $2^k$, und es ergibt sich durch Addition aller so gebildeter Gleichungen
\weq{15}{2^k\,\mathrm{sn}\,\Omega = \Sigma\sqrt{A}.}

Es ist noch zu bemerken, daß von den Größen $\psi$, also auch von den $A$, die Hälfte
verschwindet. Denn es ist
\[
-1 \equiv cc' \ldots \ (\mathrm{mod}\ n),
\]
und wenn man also in (14) $\alpha \equiv -1$ setzt, so folgt:
\[
\psi(-\Omega) = (\pm 1)(\pm 1)\ldots\psi(\Omega);
\]
andererseits ist aber
\[
\psi(-\Omega) = -\psi(\Omega),
\]
und folglich verschwindet $\psi(\Omega)$ immer dann, wenn unter den in (14)
vorkommenden $k$ Größen $\pm 1$ eine gerade Anzahl von negativen Einheiten enthalten
ist.

Die Wurzeln $x_{\mu,\mu'}$ können also linear durch $2^{k-1}$ Quadratwurzeln
ausgedrückt werden.

\gesperrt{Wenn $n$ eine Primzahl oder eine Potenz einer Primzahl ist, so sind die Quadrate der
Wurzeln der Teilungsgleichung rational durch die Wurzeln der Transformationsgleichung
ausdrückbar}\footnote{Vgl. über diesen Satz \gesperrt{Kronecker}, Monatsbericht der Berliner
Akademie, 19.\ Juli 1875. \gesperrt{Kronecker} macht dort auf ein Versehen aufmerksam, das sich
in einer diesen Gegenstand betreffenden Abhandlung von \gesperrt{Jacobi} (Crelle, Bd.\ 47 und
Bd.\ 50) findet.}.

Ist $n$ eine zusammengesetzte Zahl, so läßt sich die Einteilung der Wurzeln
$x_{\mu,\mu'}$ in Reihen noch weiter treiben. Ist $p$ eine in
%CONT>>
